% 讨论部分
\section{讨论与分析}

\begin{frame}{讨论：主要发现}
    \begin{enumerate}
        \item \textbf{\color{primaryblue}表格型 Q-Learning 性能与 DQN 相当}
        \begin{itemize}
            \item 准确率：62.68\% vs 62.70\% ($\Delta = -0.02\%$, $p = 0.981$)
            \item 在 9 名被试、5 次重复实验的大规模评估中得到验证
        \end{itemize}

        \vspace{0.5em}
        \item \textbf{\color{primaryblue}表格型 Q-Learning 稳定性更优}
        \begin{itemize}
            \item 标准差：8.20\% vs 9.21\%
            \item 跨被试稳定性相对提升约 11\%
        \end{itemize}

        \vspace{0.5em}
        \item \textbf{\color{primaryblue}轻量化优势显著}
        \begin{itemize}
            \item 无需训练深度神经网络
            \item 模型参数量减少 92 倍，推理延迟降低 1000 倍
        \end{itemize}

        \vspace{0.5em}
        \item \textbf{\color{primaryblue}Standard Q-Learning 优于复杂变体}
        \begin{itemize}
            \item Double Q-Learning 和 Dueling 架构未带来显著收益
            \item 任务状态空间相对简单（136 状态），标准方法已足够
        \end{itemize}
    \end{enumerate}
\end{frame}

\begin{frame}{讨论：理论与实践意义}
    \begin{columns}[T]
        \begin{column}{0.48\textwidth}
            \begin{block}{\normalsize {理论意义}}
                \begin{itemize}
                    \item \textbf{挑战深度强化学习的"必要性"假设}
                    \begin{itemize}
                        \item 并非所有任务都需要深度模型
                        \item 简单方法在特定场景下可能是更优选择
                    \end{itemize}

                    \item \textbf{提出 BCI 轻量化设计框架}
                    \begin{itemize}
                        \item 数据稀缺性约束
                        \item 设备成本约束
                        \item 辅助任务定位约束
                    \end{itemize}

                    \item \textbf{界定表格型与深度方法的适用边界}
                    \begin{itemize}
                        \item 状态空间规模 $< 10^4$：表格型更优
                        \item 状态空间规模 $> 10^6$：深度方法更优
                    \end{itemize}
                \end{itemize}
            \end{block}
        \end{column}

        \begin{column}{0.48\textwidth}
            \begin{block}{\normalsize {实践意义}}
                \begin{itemize}
                    \item \textbf{适合资源受限的嵌入式 BCI 设备}
                    \begin{itemize}
                        \item 内存占用 $< 10$ KB
                        \item 推理延迟 $< 1$ ms
                        \item 无需专用加速硬件
                    \end{itemize}

                    \vspace{0.5em}
                    \item \textbf{降低系统部署成本}
                    \begin{itemize}
                        \item 无需深度学习框架
                        \item 可轻松移植至嵌入式平台
                    \end{itemize}

                    \vspace{0.5em}
                    \item \textbf{提升在线 BCI 系统响应性}
                    \begin{itemize}
                        \item 快速个体校准
                        \item 低延迟推理
                    \end{itemize}
                \end{itemize}
            \end{block}
        \end{column}
    \end{columns}
\end{frame}

\begin{frame}{讨论：局限性与未来工作}
    \begin{columns}[T]
        \begin{column}{0.48\textwidth}
            \begin{alertblock}{\normalsize 局限性}
                \begin{itemize}
                    \item \textbf{状态空间离散化}
                    \begin{itemize}
                        \item 需要将连续时间窗边界离散化
                        \item 当前步长 (0.2s) 已足够精细
                    \end{itemize}

                    \vspace{0.5em}
                    \item \textbf{仅使用离线数据验证}
                    \begin{itemize}
                        \item 当前阶段以验证方法有效性为主
                        \item 未来需探索在线自适应优化
                    \end{itemize}

                    \vspace{0.5em}
                    \item \textbf{低表现被试性能提升有限}
                    \begin{itemize}
                        \item 这是 BCI 领域的共性问题
                        \item 需结合迁移学习等其他技术
                    \end{itemize}
                \end{itemize}
            \end{alertblock}
        \end{column}

        \begin{column}{0.48\textwidth}
            \begin{exampleblock}{\normalsize 未来工作}
                \begin{itemize}
                    \item 探索连续状态空间方法
                    \item 实现在线自适应时间窗优化
                    \item 结合迁移学习提升低表现被试性能
                    \item 在嵌入式设备上部署验证
                    \item 扩展至多模态 BCI 任务
                \end{itemize}
            \end{exampleblock}
        \end{column}
    \end{columns}
\end{frame}

\begin{frame}{研究结论}
    \begin{enumerate}
        \item \textbf{\color{primaryblue}方法创新}
        \begin{itemize}
            \item 提出基于表格型 Q-Learning 的运动想象时间窗优化方法
            \item 将时间窗选择建模为 MDP，使用 Q-Learning 自适应优化
        \end{itemize}

        \vspace{0.5em}
        \item \textbf{\color{primaryblue}性能验证}
        \begin{itemize}
            \item 与 DQN 准确率差异仅 -0.02\% ($p = 0.981$)
            \item 跨被试稳定性提升 11\%
            \item 自适应时间窗显著优于固定时间窗（提升约 37\%-41\%）
        \end{itemize}

        \vspace{0.5em}
        \item \textbf{\color{primaryblue}轻量化优势}
        \begin{itemize}
            \item 参数量减少 92 倍，推理延迟降低 1000 倍
            \item 内存占用减少 125 倍，训练时间缩短 3 倍
        \end{itemize}

        \vspace{0.5em}
        \item \textbf{\color{primaryblue}理论贡献}
        \begin{itemize}
            \item 挑战深度强化学习的"必要性"假设
            \item 提出 BCI 轻量化设计框架
            \item 界定表格型与深度方法的适用边界
        \end{itemize}
    \end{enumerate}
\end{frame}
